
\section{预备事实}

由\mbox{必要练习~15.1.C}，四个变量的三次\trForm{}构成一个 $20$ 维\trVectorSpace，于是 $\mathbb{P}^{3}$ 中的三次曲面由 $\mathbb{P}^{19}$ 参数化。设 $\Delta\subset\mathbb{P}^{19}$ 为对应于奇异三次\trHypersurface{}的\trLocus。\mbox{\trBertini{}定理~13.4.2} 表明：光滑三次曲面所在的\trLocus{} $\mathbb{P}^{19}\setminus\Delta$ 在全体三次曲面构成的空间 $\mathbb{P}^{19}$ 中既\trDense{}又\trOpen。

\textbf{27.1.A.（练习）}\quad 证明：任何光滑三次曲面 $X$ 都是「反典范嵌入的」——也就是说，它由完全反\trCanonical\trLinearSystem{} $\mathscr{K}_X^{\vee}$ 所嵌入。提示：\trAdjunctionFormula（\mbox{练习~21.5.B}）。

\textbf{27.1.B.（练习）}\quad 设 $X\subset\mathbb{P}^{3}_{\bar{k}}$ 是一张光滑三次曲面，而 $C$ 是 $X$ 上的一条\trCurve。证明：$C$ 是一条直线，当且仅当 $C$ 是一条 $(-1)$-\trCurve（\mbox{定义~20.2.8}）。提示：仍是\trAdjunctionFormula，或许以\mbox{练习~20.2.B}(a) 的面貌出现；另外还可参\mbox{练习~18.6.K}。

找出\emph{单独一个}域上\emph{单独一张}恰好具有 27 条直线的三次曲面，将会是有用的：

\textbf{27.1.C.（动手练习）}\quad 证明\trFermat{}三次曲面
\[
\tag{27.1.0.1}\label{eq:fermat-cubic-surface}
x_0^{3}+x_1^{3}+x_2^{3}+x_3^{3}=0
\]
在 $\mathbb{P}^{3}_{\mathbb{C}}$ 中恰有 27 条直线，每条都形如
\[
x_0+\omega x_i=x_j+\omega' x_k=0,
\]
其中 $\{1,2,3\}=\{i,j,k\}$，$j<k$，而 $\omega$ 与 $\omega'$ 都是 $1$ 的立方根（可以相同）。这将需要一些蛮力。可能的提示：在相差一个坐标置换的意义下，证明 $\mathbb{P}^{3}$ 中的每条直线都可以写成 $x_0=ax_2+bx_3$，$x_1=cx_2+dx_3$。证明：这条直线落在 \eqref{eq:fermat-cubic-surface} 上，当且仅当
\[
\tag{27.1.0.2}\label{eq:fermat-line-condition}
a^{3}+c^{3}+1=b^{3}+d^{3}+1=a^{2}b+c^{2}d=ab^{2}+cd^{2}=0
\]
再证明：若 $a$、$b$、$c$、$d$ 全都非零，则 \eqref{eq:fermat-line-condition} 无解。

\textbf{27.1.1.}\ \hard\ \textbf{奇异三次曲面在全体三次曲面的参数空间 $\mathbb{P}^{19}$ 中给出一个\trIrreducible\trDivisor{} $\Delta$。}\quad 这段讨论被标上星号，不是因为它难，而是因为本章其余部分并不需要它。尽管如此，它是对我们已学内容的一个漂亮应用，并且预示了 \S27.2 中论证的关键部分。

\mbox{\trBertini{}定理~13.4.2} 表明：全体三次曲面构成的空间 $\mathbb{P}^{19}$ 中光滑三次曲面所在的\trLocus{}既\trDense{}又\trOpen，但该证明还暗示了更多的东西：这个\trLocus{}的\trComplement{} $\Delta$ 具有纯余维数 $1$。我们现在就来证明事情确乎如此，甚至证明这个「\trDiscriminant{}超曲面」$\Delta$ 是\trIrreducible{}的。

\textbf{27.1.D.（练习）}\quad（两小问共用提示：回想\mbox{练习~12.2.N} 的解答。）
（a）定义\trIncidenceCorrespondence{} $Y\subset\mathbb{P}^{19}\times\mathbb{P}^{3}$，它对应于「一张三次曲面 $X$ 连同其上一个\trSingularPoint{} $p\in X$」这份数据。（这是练习的一部分！我们需要的是作为概形的 $Y$，而不只是作为集合的 $Y$。）令 $\mu$ 为\trProjection{} $Y\to\mathbb{P}^{19}$。
\[
\begin{tikzcd}[row sep=2.2em, column sep=3.6em]
& Y \ar[dl, "\mu"] \ar[dr, "\mathbb{P}^{15}\text{-丛}"] & \\
\mathbb{P}^{19} && \mathbb{P}^{3}
\end{tikzcd}
\]
（b）通过把 $Y$ 描述为 $\mathbb{P}^{3}$ 上的一个 $\mathbb{P}^{15}$-丛，证明 $Y$ 是维数为 $18$ 的\trIrreducible{}光滑\trVariety。

\textbf{27.1.E.（练习）}\quad 证明：存在只有一个\trSingularPoint{}的三次曲面。你尽可随意假设 $k$ 是你最喜欢的\trField；这里的要点是让你熟悉几何想法，而不是正特征\trField{}的种种怪癖。（$\operatorname{char}k\neq 2$ 情形的提示：\mbox{练习~13.3.C}。）

\textbf{27.1.F.（练习）}\quad 证明：$\mu(Y)=\Delta$ 是一个余维数为 $1$ 的\trClosed\trIrreducible{}子集。余维数 $1$ 部分的提示：\mbox{练习~12.4.A} 或\mbox{定理~12.4.1}，与\mbox{练习~27.1.E} 结合起来用。

你的论证几乎无须改动即可推广，用以处理 $\mathbb{P}^{n}$ 中次数 $d$ 的\trHypersurface。
